%mac
\begin{dang}{PHƯƠNG TRÌNH VÀ ĐỒ THỊ CỦA CHUYỂN ĐỘNG THẲNG ĐỀU}
\setcounter{paragraph}{0}
\paragraph{Vị trí và thời điểm hai xe gặp nhau}
\begin{itemize}
\item Viết phương trình chuyển động của hai xe trên cùng một hệ quy chiếu
\begin{align*}\begin{cases}
x_1=x_{01}+v_1t.\\
x_2=x_{02}+v_2t.
\end{cases}
\end{align*}
\item Khi hai xe gặp nhau thì $x_1=x_2$.
\item Giải phương trình tìm t và x.
\end{itemize}
\paragraph{Phương pháp vẽ đồ thị}
\begin{itemize}
\item Vẽ hệ trục tọa độ $Odt$.
\item Lập bảng độ dịch chuyển - thời gian $(d,\ t)$. 
\item Vẽ đồ thị tọa độ bằng cách vẽ đường thẳng hoặc các đoạn thẳng, nửa đường thẳng qua từng cặp điểm đã xác định.
\end{itemize}
\begin{note}
\begin{itemize}
\item[$\bullet$] Đồ thị độ dịch chuyển - thời gian của chuyển động thẳng đều là đường thẳng, do đó ta chỉ cần xác định hai điểm trên đường thẳng đó là đủ.
\item[$\bullet$] Trừ trường hợp đặc biệt trong quá trình chuyển động vật dừng lại một thời gian hoặc thay đổi tốc độ, khi đó ta phải xác định các cặp điểm khác.
\end{itemize}
\end{note}
\begin{center}
\begin{tabular}{|c|c|c|c|}
\hline 
\textbf{Đồ thị dốc lên} & \textbf{Đồ thị dốc xuống} & \textbf{Đồ thị nằm ngang} & \textbf{Đồ thị song song}\\ 
\hline 
$v > 0$ & $v < 0$ & $v = 0$ & $v=const$\\ 
\hline 
Vật đi theo $\oplus$  & Vật đi theo $\circleddash$ & Vật đứng yên & Các vật có cùng vận tốc\\ 
\hline 
\end{tabular} 
\end{center}
\end{dang}